% =====================================================================
% RASCUNHO -- esqueleto (backbone) da fundamentação teórica
% ---------------------------------------------------------------------
% Dois capítulos:
%   1. MIMO massivo e equalização
%   2. Equalização descentralizada sob ruído colorido
% Os outros capítulos (introdução, resultados etc.) ficam como estão.
% Nenhum texto de corpo novo foi escrito aqui: só a estrutura e, em
% comentário, o que cada seção deve conter e de onde vem o texto existente.
%
% Base:
%   - capítulos atuais: 1-noise-and-constraints.tex ("coisas") e
%     0.channel-modelling.tex;
%   - roteiro T1-T6 de constraints/notas_redacao.md (ruído colorido até G_c);
%   - proposta do TG1 (TG1_Português_ECP_UFRGS/.../sbc-template.tex),
%     "Matrix Inversion Techniques".
%
% Legenda dos comentários de cada seção:
%   CONTEÚDO  o que a seção precisa dizer
%   TEXTO     de onde vem o texto já escrito (arquivo e seção atual)
%   ESTADO    escrito / parcial / vazio / a escrever
%
% Rótulos: os que já existem nos capítulos atuais foram mantidos
% (sec:equalization, subsec:bloco-coerencia, subsec:construcao-R,
% subsec:o-que-R-representa) e sec:sysmodel foi criado (hoje é citado e não
% existe). Por isso este arquivo substitui 1-noise-and-constraints.tex e
% 0.channel-modelling.tex; não incluir os três juntos.
% =====================================================================


% =====================================================================
\chapter{MIMO massivo e equalização}\label{chap:mimo-equalizacao}
% =====================================================================
% Base geral, válida para qualquer receptor MIMO: o modelo do sistema, o canal
% e os equalizadores lineares. Termina no LMMSE e na SER, que o capítulo
% seguinte usa.
% Ordem: cada seção usa só o que veio antes.

% ---------------------------------------------------------------------
\section{Modelo de sistema}\label{sec:sysmodel}
% ---------------------------------------------------------------------
% CONTEÚDO  Uplink multiusuário: y = Hs + n; dimensões M, K; símbolos de uma
%           constelação (16-QAM, QPSK); E_s; SNR = 10 log10(E_s / sigma^2).
% TEXTO     1-noise-and-constraints.tex, "Uplink/Downlink" (11 linhas).
% ESTADO    parcial. "SNR" não aparece hoje no capítulo.

\subsection{Enlace de subida}
\subsection{Símbolos e constelação}
\subsection{Relação sinal-ruído}

% ---------------------------------------------------------------------
\section{Modelo de canal}\label{sec:modelo-canal}
% ---------------------------------------------------------------------
% CONTEÚDO  O que o canal representa; modelos determinísticos e
%           estocásticos; desvanecimento; bloco de coerência; arranjos de
%           antenas; modelo 3GPP, QuaDRiGa e cenários.
% TEXTO     0.channel-modelling.tex (inteiro);
%           1-noise-and-constraints.tex, "Multipath, desvanecimento..." e
%           "Bloco de Coerência" (vazias);
%           constraints/rascunho_bloco_canal.tex (partes 1 e 2; ajustar o tom).
% ESTADO    parcial

\subsection{Tipos de modelo de canal}
\subsection{Desvanecimento}
\subsection{Bloco de coerência}\label{subsec:bloco-coerencia}
% O setup de simulação se apoia aqui: W é reaproveitado para vários símbolos.
\subsection{Arranjos de antenas}
\subsection{Modelo 3GPP e simulador QuaDRiGa}
\subsection{Cenários de propagação}
% A frase "este trabalho usa o cenário UMa NLOS" vai para o setup.

% ---------------------------------------------------------------------
\section{Equalização e detecção}\label{sec:equalization}
% ---------------------------------------------------------------------
% CONTEÚDO  Detector ML e por que é inviável; separação entre equalização e
%           detecção; filtros lineares; detecção pelo ponto mais próximo; SER.
% TEXTO     1-noise-and-constraints.tex, "Equalização" e "Filtros Lineares".
% ESTADO    parcial (ZF, Wiener e LMMSE escritos; SER não definida)

\subsection{Detector de máxima verossimilhança}
\subsection{Zero-forcing}
\subsection{Covariância}
\subsection{Filtro de Wiener}
\subsection{Equalizador LMMSE}
\subsection{Detecção e taxa de erro de símbolo}
% Novo: decisão pelo ponto mais próximo da constelação; definição de SER.

% =====================================================================
\chapter{Equalização descentralizada sob ruído colorido}\label{chap:dbp-ruido}
% =====================================================================
% Núcleo do trabalho: o ruído colorido, a estimação da sua covariância, a
% arquitetura descentralizada, o BCD-MMSE e a matriz G_c que cada DU inverte,
% e as técnicas para inverter essa matriz.
% Roteiro T1-T6 de constraints/notas_redacao.md.


% ---------------------------------------------------------------------
\section{Ruído colorido}\label{sec:ruido-colorido}
% ---------------------------------------------------------------------
% CONTEÚDO  Do AWGN ao ruído colorido: interferentes + ruído térmico;
%           construção de R = beta Hbar Hbar^H + sigma^2 I; IoT; leitura de R.
% TEXTO     1-noise-and-constraints.tex, "Modelo AWGN" e "Modelo do ruído
%           colorido" (assinatura espacial, hipóteses, construção de R_n).
% ESTADO    escrito até a construção de R; falta a IoT

\subsection{Ruído branco (AWGN)}
\subsection{Assinatura espacial e interferência}
\subsection{Hipóteses do modelo}
\subsection{Construção da matriz de covariância}\label{subsec:construcao-R}
\subsection{Interferência sobre ruído térmico (IoT)}
% Novo: definição da IoT; papel de beta.
\subsection{O que a matriz de covariância representa}\label{subsec:o-que-R-representa}

% ---------------------------------------------------------------------
\section{Estimação da matriz de covariância}\label{sec:estimacao-R}
% ---------------------------------------------------------------------
% CONTEÚDO  Roteiro T1: R não pode ser calculada; estimativa amostral com N
%           amostras de piloto; posto de R-chapéu; quantidade e custo de N.
% TEXTO     1-noise-and-constraints.tex, "Estimação de R_n" (subseções
%           parciais); constraints/rascunho_filtros_covariancia.tex.
% ESTADO    parcial

\subsection{Por que estimar}
\subsection{Amostras de piloto}
% Absorve a seção "Pilotos" (vazia) do capítulo atual.
\subsection{Intervalo das amostras}
\subsection{O que cada amostra contribui}
\subsection{Quantidade de amostras}
% Posto de R-chapéu <= N: motiva o condicionamento (subseção seguinte).
\subsection{Condicionamento da estimativa}\label{subsec:condicionamento-R}
% CONTEÚDO  Quando N < M, o posto de R-chapéu é no máximo N e a matriz não
%           tem inversa (M = 256, N = 192). O LMMSE precisa dessa inversa.
%           Shrinkage de Ledoit-Wolf: R-chapéu vira (1 - rho) R-chapéu +
%           rho (tr R-chapéu / M) I, com rho calculado a partir das próprias
%           amostras (sem constante a escolher). Alternativa: carga diagonal
%           fixa, que depende de um epsilon.
% TEXTO     nenhum na fundamentação; parágrafo sobre Ledoit-Wolf em
%           rascunho_setup_simulacao.tex (subseção Receptor);
%           dbp_sim/params.py, M6_Covariance.shrinkage.
% FONTE     Ledoit e Wolf (2004), J. Multivariate Analysis 88(2) -- falta
%           no .bib; posto <= N: livro-texto de álgebra linear.
% MÓDULO    M6 (Conditioner em dbp_sim/m5_noise.py)
% ESTADO    a escrever
\subsection{Custo das amostras}

% ---------------------------------------------------------------------
\section{Processamento descentralizado}\label{sec:dbp}
% ---------------------------------------------------------------------
% CONTEÚDO  Arquitetura DBP: clusters, DUs, CU; estrela e daisy chain;
%           roteiro T2 (o que cada DU alcança) e T3 (a bifurcação: descartar
%           os blocos cruzados -> BDAC; mantê-los comprimidos -> BCD).
% TEXTO     1-noise-and-constraints.tex, "Definição Decentralizada" (vazia)
%           e "Estimação nas DUs" (2 linhas).
% ESTADO    vazio

\subsection{Arquitetura e particionamento das antenas}
\subsection{Estimação da covariância nas DUs}
\subsection{Aproximação bloco-diagonal (BDAC-MMSE)}

% ---------------------------------------------------------------------
\section{Equalizador BCD-MMSE}\label{sec:bcd}
% ---------------------------------------------------------------------
% CONTEÚDO  Roteiro T4 (formulação BCD), T5 (da eq. 26 à eq. 30) e T6 (o
%           fator a inverter, G_c = E_s H_c H_c^H + R_cc, e suas
%           propriedades: hermitiana, definida positiva).
% TEXTO     constraints/notas_redacao.md (T4-T6);
%           constraints/modelo_ruido_e_hipoteses.md, seção 10.
% ESTADO    a escrever

\subsection{Formulação por blocos}
\subsection{Atualização de cada DU}
\subsection{A matriz invertida em cada DU}\label{subsec:Gc}

% ---------------------------------------------------------------------
\section{Inversão de matrizes}\label{sec:inversao}
% ---------------------------------------------------------------------
% CONTEÚDO  Custo O(M_c^3); decomposição de Cholesky para matrizes
%           hermitianas definidas positivas; substituições direta e reversa;
%           série de Neumann e outras alternativas; por que Cholesky para G_c.
% TEXTO     Proposta do TG1, "Matrix Inversion Techniques".
% ESTADO    a escrever (traduzir e adaptar)

\subsection{Decomposição de Cholesky}
\subsection{Série de Neumann}
\subsection{Comparação}
